\documentclass[11pt,a4paper]{article}
\usepackage[utf8]{inputenc}
\usepackage[english]{babel}
\usepackage{amsmath}
\usepackage{amsfonts}
\usepackage{graphicx}
\usepackage{geometry}
\usepackage{listings}
\usepackage{booktabs}

\geometry{margin=1in}

\title{\textbf{The Circular Compartment Train (CCT):}\\A Novel Sub-Linear Geometric Data Structure for Unique Set Topologies}
\author{\textbf{Nils Gischler} \\ Theoretical Computer Science Division}
\date{\today}

\begin{document}

\maketitle

\begin{abstract}
This paper presents the \textit{Circular Compartment Train (CCT)}, a highly adaptive, sub-linear data structure tailored for dynamic numeric sets. By blending a cyclically double-linked pointer system (the ``train'') with a dynamic centralized geometric partitioning mechanism (the ``compartments''), the CCT solves the linear lookup vulnerabilities inherent in standard linked layouts. I present two design topologies: the \textit{CCT-Bag} for high-frequency multiset data streams and the \textit{CCT-Set}, which leverages an immediate $\mathcal{O}(1)$ bitmap exclusion layout to prevent structural duplicate inflation. I prove that the balance bound governed by $M = \sqrt{N} + (0.05 \cdot N)$ preserves stable $\mathcal{O}(\sqrt{N})$ boundaries via structural compartment splits and flexible boundary shifts.
\end{abstract}

\section{Introduction}
Linear scanning limits standard linked list variants to an unfavorable $\mathcal{O}(N)$ lookup boundary. While structures such as Skip Lists or self-balancing Binary Search Trees (BST) bound operations to logarithmic intervals, they suffer from considerable structural pointer ballast and strict top-down layout constraints. We propose the Circular Compartment Train (CCT), an expressive hybrid topology that frames data elements into a rolling, continuous railroad ring, managed natively via an analytical multi-sector coordinate index representing radial pie slices mapping value intervals.

\section{Structural Topology \& Component Architecture}
The data structure partitions architectural responsibilities across two structural segments:
\begin{enumerate}
    \item \textbf{The Centralized Geometric Index:} A dynamic registry that partitions the total numeric domain into distinct intervals called \textit{Compartments}. Each compartment maintains bounds $[\text{low}, \text{high}]$, a size tracker, and a static entry pointer targeting the exact boundary element in the node ring.
    \item \textbf{The Cyclic Train System:} A continuous, closed loop composed of doubly linked elements (\textit{Wagons}), natively permitting bidirectionality ($\text{next}$ and $\text{prev}$ pointers). The physical terminus of the sequence loops back seamlessly into the initial node element, eliminating hard structural boundary limits.
\end{enumerate}

\section{Mathematical Balancing Operations}
To maintain the runtime balance against targeted value clusters, individual compartment sizes are bounded by a dynamic equilibrium function $M$. The maximal wagon count per slice is tightly coupled to the square root of the system volume $N$ enriched by an adaptive elasticity cushion:
\begin{equation}
    M = \sqrt{N} + (0.05 \cdot N)
\end{equation}

When an ingestion node hits the threshold limit where $|C_i| > M$, the CCT framework initializes a structural splitting routine:
\begin{equation}
    \text{Threshold}_{\text{mid}} = \left\lfloor \frac{\text{low} + \text{high}}{2} \right\rfloor
\end{equation}
The bounds are recalculated instantly, shifting the boundary definitions along the ring without forcing physical element allocation moves in memory.

\section{Ingestion Filtering \& Duplicate Invalidation}
The \textit{CCT-Set} model optimizes membership operations via a bit-vector signature mask assigned to each radial compartment slice. Before traversing the cyclic node network, incoming items undergo immediate indexing lookup:
\begin{equation}
    f(X) = \text{Bitmap}[X \pmod K]
\end{equation}
If the coordinate filter returns an active state, the system flags a collisions threat and rejects the element instantly. This short-circuit strategy bounds duplicate rejection tasks to a deterministic time complexity of $\mathcal{O}(1)$.

\section{Complexity Overview}
A systematic comparison underlines the competitive profile of the Circular Compartment Train against classical layouts under uniform distributions:

\begin{table}[h]
\centering
\caption{Algorithmic Comparison Summary}
\begin{tabular}{llll}
\toprule
\textbf{Operation} & \textbf{Standard Linked List} & \textbf{Balanced BST} & \textbf{CCT (Ours)} \\
\midrule
Search Time & $\mathcal{O}(N)$ & $\mathcal{O}(\log N)$ & $\mathcal{O}(\sqrt{N})$ \\
Insertion Time & $\mathcal{O}(1)$ or $\mathcal{O}(N)$ & $\mathcal{O}(\log N)$ & $\mathcal{O}(\sqrt{N})$ \\
Duplicate Rejection & $\mathcal{O}(N)$ & $\mathcal{O}(\log N)$ & $\mathcal{O}(1)$ \\
Memory Overlap Space & Moderate & High & Low-to-Moderate \\
\bottomrule
\end{tabular}
\end{table}

\section{Conclusion}
The CCT architecture offers a compelling paradigm shift in layout design, blending cyclic connectivity with radial space splitting. Future milestones include the evaluation of dynamic slice merges under extreme deletion load states.

\end{document}
